\specialchapter{Prologue}{What it means to dissolve a problem}{4}
\begin{refsection}[chapters/00-prologue.bib]

Some problems are hard because they are hard. Others are hard because they are being asked in the wrong place. This book is about the second kind, and about a skill that research mathematicians use constantly but almost never teach: finding the setting in which a problem stops being difficult.

\addsec{Two ways to open a nut}

In \emph{R\'ecoltes et Semailles}, the long memoir he wrote in the mid-1980s, Alexander Grothendieck described two ways of opening a hard nut \parencite{Grothendieck2022}.\sidenote{The memoir circulated privately for decades and was published by Gallimard only in 2022. For an English account of these passages and of what they meant for Grothendieck's mathematics, see \textcite{McLarty2007}.} The first is the familiar one: set a chisel against the shell and strike, and if the shell holds, strike again at another point until it cracks. The second is to put the nut in a softening liquid and wait. Over weeks and months the shell softens, and eventually it opens under the pressure of a hand. Grothendieck preferred a larger image of the same kind. A sea rises so slowly that nothing seems to happen, until the problem is surrounded and submerged, dissolved in a theory that reaches far beyond what was originally to be proved.

Both methods are honorable, and much of the best mathematics ever written was cracked rather than soaked. But the two methods leave different things behind. A cracked problem leaves a proof. A dissolved problem leaves a \emph{setting}: a definition, an object, a point of view, in which the original problem and many of its neighbors have short proofs. The setting can be carried to the next problem. The chisel marks cannot.

Settings are rarely taught. Textbooks present mathematics after the sea has risen, so the reader sees the dissolved theorem and the theory that dissolved it, but not the shore that was there before. Histories present the shore but rarely isolate the move that flooded it. This book does both, one case at a time. Each chapter takes a problem that was fought with hammer and chisel, shows the fight, isolates the single reframing that ended it, re-proves the old results in the new setting, and then asks what the reframing cost and where it fails. The reframings are called \term{moves}, and the goal is that you learn to make them.

\addsec{A first case: the needle and the noodle}

The smallest example of the whole pattern fits on two pages.

\paragraph{Before.} In 1777 Georges-Louis Leclerc, comte de Buffon, published the following problem \parencite{Buffon1777}.\sidenote{Buffon had considered a related game, \emph{franc-carreau}, in a note of 1733; the needle problem appears in his \emph{Essai d'arithm\'etique morale}, a supplement to his \emph{Histoire naturelle}.} The plane is ruled by parallel lines at distance $d$. A needle of length $\ell\le d$ is dropped at random. What is the probability that it crosses a line?

Make ``at random'' precise. Let $y$ be the distance from the midpoint of the needle to the nearest line, and let $\theta$ be the angle between the needle and the lines; take $y$ uniform on $[0,d/2]$ and $\theta$ uniform on $[0,\pi/2]$, independently. The needle crosses a line exactly when $y\le\frac{\ell}{2}\sin\theta$, so
\[
  \Pr[\text{cross}] = \frac{2}{\pi}\int_0^{\pi/2}\frac{2}{d}\cdot\frac{\ell}{2}\sin\theta\dd\theta = \frac{2\ell}{\pi d}.
\]
The computation is short and correct, and it explains nothing. Why does $\pi$ appear? Why is the answer linear in $\ell$? Bend the needle slightly and the integral changes into a different integral, which no longer has a closed form.

\paragraph{The move.} In 1860 Joseph-\'Emile Barbier asked a more general question \parencite{Barbier1860}. Drop a rigid wire of arbitrary shape and length $L$, a \term{noodle},\index{Buffon's noodle} and count its crossings $N$ with the lines. What is $\E N$? The question sounds harder than Buffon's, because it is about every shape at once. It is easier, because a question about every shape at once has more structure: the answer has to be compatible with cutting noodles into pieces and gluing pieces into noodles. P\'olya called this phenomenon the \emph{inventor's paradox}: the more ambitious problem can be the easier one \parencite{Polya1945}.

\paragraph{After.} Drop a noodle by choosing a uniformly random rotation and an independent, uniformly random translation perpendicular to the lines (modulo $d$). A \emph{polygonal noodle} is a noodle made of finitely many straight edges.

\begin{theorem*}[Barbier]
Every polygonal noodle of total length $L$ satisfies $\E N = 2L/(\pi d)$.
\end{theorem*}

\begin{proof}
The plan is to show that $\E N$ is proportional to length, and then to find the constant from a single shape whose crossing number we know. Write $f(\ell)$ for the expected number of crossings of a straight segment of length $\ell$. Each edge of a randomly dropped polygonal noodle is, on its own, a randomly dropped segment, even though the edges move together; and crossings of the noodle are crossings of its edges, apart from an event of probability zero. Linearity of expectation does not care about the dependence between the edges, so
\[
  \E N = \sum_{\text{edges } e} f(\ell_e).
\]
Cutting a segment into two collinear pieces gives $f(\ell_1+\ell_2) = f(\ell_1)+f(\ell_2)$, and a longer segment contains a shorter one, so $f$ is nondecreasing. An additive, nondecreasing function on $[0,\infty)$ is linear (\cref{ex:prologue:cauchy}), so $f(\ell) = c\ell$ for a constant $c$, and therefore $\E N = cL$ for every polygonal noodle.

To find $c$, look at a circle $C$ of diameter $d$. It is not polygonal, but it is squeezed between polygons. Let $P_n$ and $Q_n$ be the regular $n$-gons inscribed in and circumscribed about $C$. Each of the three closed convex curves meets a line in $0$ or $2$ points, apart from an event of probability zero; a line that meets $P_n$ meets $C$, and a line that meets $C$ meets $Q_n$. The circle has width exactly $d$ in every direction, so it meets exactly one line, twice. Hence
\[
  c\cdot\operatorname{length}(P_n) = \E N(P_n) \le 2 \le \E N(Q_n) = c\cdot\operatorname{length}(Q_n).
\]
Both perimeters tend to $\pi d$ as $n\to\infty$, so $c\,\pi d = 2$.
\end{proof}

Buffon's answer is now a remark. A needle with $\ell\le d$ crosses at most one line, so the probability that it crosses equals the expected number of crossings, $2\ell/(\pi d)$. The answer is linear in $\ell$ because the expectation is additive, and $\pi$ appears because the calibrating shape is a circle. Both of the questions that the integral left open now have answers.

\paragraph{Costs and limits.} The difficulty did not vanish; it moved. Part of it went into linearity of expectation, which holds for dependent pieces, and part into the existence of one shape with an obvious answer: a curve of constant width. Paying that once buys every noodle at the same price. The move also has a limit. It computes the expected number of crossings, and for a needle longer than $d$ the probability of crossing is no longer an expectation, so Buffon's long-needle problem still requires an honest integral. Every chapter of this book ends with the same two questions: \emph{where did the difficulty go?}, and \emph{what did not dissolve?}

\addsec{The dissolution test}

Not every reframing dissolves a problem. Some merely rename it. A definition that builds the conclusion into the hypotheses makes every theorem a tautology and every application as hard as before. The cases in this book were chosen because they pass a test with three gates and a dividend. A reframing must pay at least one of the following tolls.\index{dissolution test}

\begin{description}[leftmargin=1.2em,style=nextline,font=\normalfont\scshape\color{sea}]
\item[Shorter proofs.] Old theorems receive proofs that are shorter for a reason, not merely in notation. Arzel\`a's convergence theorem for Riemann integrals has elementary proofs that are long and delicate; after Lebesgue it is a two-line corollary of the dominated convergence theorem (\cref{ch:lebesgue}). Barbier's argument replaces an integral by linearity.
\item[Fewer assumptions.] Hypotheses that looked necessary turn out not to be. Uniform convergence is replaced by domination; the straightness of the needle is not used at all.
\item[Unification.] Results that were proved separately become one statement. K\"onig's theorem on bipartite matchings, the max-flow min-cut theorem, Dilworth's theorem and von Neumann's minimax theorem are four faces of linear programming duality (\cref{ch:lpduality}).
\end{description}

\noindent The dividend is \textsc{new reach}: results that the old setting could not even state. Barbier's formula, read backwards, measures the length of every rectifiable curve by counting how often it crosses random lines; Lebesgue's integral makes $L^2$ complete and Fourier series a statement about a Hilbert space.

Grothendieck added an observation that the cases bear out. The decisive idea is often not grand but humble; he recalled that the idea of a scheme was so humble that nobody before him had stooped low enough to pick it up, and that for years it did not look serious \parencite{Grothendieck2022}. The move is to \emph{stoop low}: to pick up the obvious object that nobody bothered to pick up. A needle becomes a noodle. A graph becomes a function on the unit square (\cref{ch:graphon}). An adversary becomes a binary tree of labelled points (\cref{ch:littlestone}). A dataset becomes a list whose order does not matter (\cref{ch:conformal}). A set becomes the zero set of a polynomial (\cref{ch:dimension}).

\addsec{Where did the difficulty go?}

Difficulty is conserved. When a hard theorem acquires a short proof, the work it used to require has gone somewhere, and the honest question is where. In a genuine dissolution the difficulty is paid once, in building the setting, and is then amortized over many problems. Constructing Lebesgue measure is real work, but it is done once, and then it pays for Fourier analysis, probability theory and the calculus of variations. The Borsuk--Ulam theorem is real topology, but once it is available Kneser's conjecture, the necklace-splitting theorem and the ham sandwich theorem cost a page each (\cref{ch:kneser}).

A fake dissolution charges the full price again at every use. If verifying the hypotheses of a general theorem is as hard as proving the special case directly, the general theorem has relocated the difficulty, not dissolved it. This is the most common way abstraction fails, and the cases in this book are chosen, in part, to show what the alternative looks like.

Two further warnings. A generalization whose proof is longer than the special case is decoration, not dissolution. And a nut that is already open should not be soaked: a clean, short, direct proof does not need a theory, and burying it in one makes mathematics worse. The \emph{Costs and limits} section of each chapter is there to keep the book honest about both.

\addsec{The anatomy of a case}

Every chapter after this one is a case with the same anatomy, so that the moves can be compared.

\begin{description}[leftmargin=1.2em,style=nextline,font=\normalfont\scshape\color{sea}]
\item[Before.] The problem in its original setting, with enough of the old mathematics to feel its resistance: old statements, an old proof or its skeleton, the places where it is ad hoc, and the questions it could not answer.
\item[The move.] The reframing: a new object, definition or viewpoint, and why it is natural in hindsight.
\item[After.] The old results re-proved, visibly shorter, followed by results the old setting could not reach.
\item[Costs and limits.] Where the difficulty went, what did not dissolve, and which other frameworks compete.
\item[The move, distilled.] The move as a transferable recipe, in a box like the one below, followed by the chapter's dissolution test.
\item[Exercises.] Three kinds: \textsc{consolidation} (re-derive results of the chapter in the new setting), \textsc{transfer} (apply the same move to a new problem, which is the point of the book), and \textsc{research} (open-ended questions, some of them open problems). Stars mark difficulty.
\item[Notes and References.] History and pointers to the literature.
\end{description}

\noindent Here is the move of this prologue, distilled in the format used throughout.

\begin{move}{Generalize to simplify}
\movepart{In one sentence} Replace a question about one object by the same question about a family with more structure, and let the structure force the answer.
\movepart{Triggers} A computation that produces the right number but no reason for it; an answer that is suspiciously linear or suspiciously symmetric; a special case that breaks as soon as the object is perturbed.
\movepart{Recipe} Enlarge the class of objects until the quantity you want becomes additive, invariant or otherwise constrained; find the constraint; calibrate with one object whose answer is obvious.
\movepart{Anti-pattern} A generalization chosen for coverage rather than structure, whose proof is longer than the special case.
\end{move}

\dissolutiontest{An integral over angles becomes a two-step argument: additivity, then calibration.}{The needle no longer needs to be straight, or even polygonal after an approximation.}{Needles, polygons and circles obey one formula.}{Expected crossings of every rectifiable curve, and hence a way to measure length by counting crossings.}

\addsec{Four families of moves}

The seventeen cases are grouped into four parts according to the kind of move they teach. The groups are not rigid. Most dissolutions use more than one move, and the epilogue returns to how the moves combine.

\begin{fullwidth}
\small
\renewcommand{\arraystretch}{1.18}
\begin{tabularx}{\linewidth}{@{}r>{\raggedright}p{0.36\linewidth}>{\raggedright}p{0.2\linewidth}X@{}}
\toprule
 & \textsc{case} & \textsc{move} & \textsc{field}\\
\midrule
\multicolumn{4}{@{}l}{\textsc{\color{sea}Part I. Building the right object}}\\
1 & Lebesgue's integral & completion & analysis\\
2 & Graph limits & completion & combinatorics\\
3 & Kolmogorov's axioms & axiomatization & probability\\
4 & Matroids & axiomatization & algorithms\\
\multicolumn{4}{@{}l}{\textsc{\color{sea}Part II. Changing the battlefield}}\\
5 & Linear programming duality & duality & optimization, TCS\\
6 & No-regret learning & existence to dynamics & learning, games\\
7 & Kneser's conjecture & transfer to topology & combinatorics\\
8 & Boolean Fourier analysis & diagonalize the symmetry & TCS, social choice\\
9 & Oddtown to cap sets & linearization & combinatorics\\
10 & Optimal transport & geometrization & analysis, probability\\
\multicolumn{4}{@{}l}{\textsc{\color{sea}Part III. Finding the right measure}}\\
11 & VC dimension & the right invariant & learning theory\\
12 & Littlestone dimension & the right invariant & learning, privacy, logic\\
13 & Le Cam's experiments & localize to a limit & statistics\\
14 & Conformal prediction & exploit a free symmetry & statistics, ML\\
\multicolumn{4}{@{}l}{\textsc{\color{sea}Part IV. From existence to construction}}\\
15 & Entropy compression & encoding & algorithms, probability\\
16 & Sum of squares & proofs to algorithms & TCS, statistics\\
17 & Pseudorandomness & unification & complexity theory\\
\bottomrule
\end{tabularx}
\end{fullwidth}

\addsec{How to read this book}

The chapters are independent, and each defines what it uses. The prerequisites are undergraduate analysis, linear algebra, probability and algorithms, together with the habit of reading proofs. A few chapters lean on earlier ones (\cref{ch:kolmogorov} uses the integral of \cref{ch:lebesgue}, and \cref{ch:regret} continues \cref{ch:lpduality}), and the margins point out the connections as they arise.

Readers with a particular field in mind may prefer a path. For learning theory and machine learning: Chapters 5, 6, 11, 12, 14 and 16, then 10 and 13. For theoretical computer science: Chapters 4, 5, 7, 8, 9, 15, 16 and 17. For statistics: Chapters 3, 11, 13, 14 and 16. For analysis and probability: Chapters 1, 2, 3, 10 and 13. Every path should end with the epilogue, which compares the moves.

Side notes in the margin carry history, cross-references and asides; nothing in them is needed to follow the mathematics. Proofs follow Halmos's advice \parencite{Halmos1970}: they say where they are going, they argue forward, and they end with a tombstone. The exercises are part of the text. The \textsc{transfer} exercises in particular are where the moves are learned, because recognizing a move in a solved case is much easier than making it in an unsolved one.

\begin{exercises}
\item\label{ex:prologue:cauchy} \consolidation Let $f\colon[0,\infty)\to[0,\infty)$ satisfy $f(a+b)=f(a)+f(b)$ for all $a,b\ge 0$, and suppose $f$ is nondecreasing. Prove that $f(x)=cx$ with $c=f(1)$. \hint{First show $f(q)=cq$ for rational $q$, then squeeze.}
\item \consolidation A wire bent into an equilateral triangle of side $s<d$ is dropped. Show that it meets a line with probability $3s/(\pi d)$. Compare with a straight needle of the same total length $3s\le d$, and explain the factor of two.
\item \consolidation A needle of length $\ell = 2d$ is dropped. Show that $\E N = 4/\pi$, and explain why this does not determine $\Pr[N\ge 1]$.
\item \transfer\hard Let $K$ be a compact convex set in the plane with perimeter $p$, and let a random line be chosen by dropping $K$ on the ruled plane with spacing $d$ larger than the diameter of $K$. Use Barbier's theorem to show that the probability that $K$ meets a line is $p/(\pi d)$. Deduce Cauchy's formula: the average width of $K$ over all directions is $p/\pi$.
\item \transfer\hard A rectifiable curve is dropped on the ruled plane. Show that its length can be estimated without bias by counting crossings: $L = \frac{\pi d}{2}\,\E N$. Extend Barbier's theorem from polygonal noodles to rectifiable curves using inscribed polygons and monotone convergence. \hint{If a chord of the curve crosses a line, so does the arc between its endpoints.}
\item \transfer Find another classical computation in which enlarging the question makes the answer additive, and redo it by the method of this prologue. One candidate: the expected number of cycles of a uniformly random permutation of $\{1,\dots,n\}$. Give each element the weight $1/k$, where $k$ is the length of its cycle, so that the number of cycles becomes a sum over elements, and show that the expectation is $1+\frac12+\dots+\frac1n$.
\item \research Choose a theorem in your own field whose proof you regard as a fight. Write down the setting in which it is stated, list the hypotheses the proof actually uses, and ask which object, if it existed, would make the theorem a special case of something easier. Keep the list; the epilogue asks you to come back to it.
\end{exercises}

\notesandrefs

Grothendieck's images of the nut and the rising sea come from \emph{R\'ecoltes et Semailles} \parencite{Grothendieck2022}; \textcite{McLarty2007} translates the passages and explains how they describe the architecture of Grothendieck's work on the Weil conjectures. Buffon's problem appeared in \textcite{Buffon1777}, and the noodle argument in \textcite{Barbier1860}; \textcite{AignerZiegler2018} give the version with inscribed and circumscribed polygons used above. P\'olya's inventor's paradox is one of the heuristics in \textcite{Polya1945}, a book whose spirit this one shares at a different level.

Two books are close relatives. \textcite{Lakatos1976} studies, through one long case, how definitions and theorems are adjusted to each other under the pressure of counterexamples; this book studies how a change of setting can make the pressure disappear. \textcite{Hadamard1945} examined how mathematicians find their ideas, from the inside. Neither is a textbook of moves, and neither attempts to make the moves transferable through exercises, which is the purpose of what follows.

\printchapterbib
\end{refsection}
